\documentclass[11pt,a4paper]{article}
\usepackage[utf8]{inputenc}
\usepackage[T1]{fontenc}
\usepackage[margin=1.6cm]{geometry}
\usepackage{xcolor}
\usepackage{enumitem}
\usepackage{hyperref}
\setlength{\parindent}{0pt}
\setlength{\parskip}{4pt}
\definecolor{primary}{HTML}{1F2937}
\definecolor{secondary}{HTML}{5B6573}
\begin{document}
{\color{primary}\LARGE\textbf{Gustavo Vega}}\\
{\color{secondary}\large Operaciones de supermercados | Pedidos, stock y mejora de procesos}\\
gusvegabol@gmail.com \textbar{} 669 549 933

\section*{Perfil}
Profesional de operaciones de supermercados con experiencia en previsión de ventas, pedidos, stock, rotación y mejora de procesos entre tiendas.
Disponibilidad para turnos rotativos de mañana o tarde.\\

\section*{Experiencia operativa relevante}
\begin{itemize}[leftmargin=*,nosep]
\item Diseñé una base de datos de ventas diarias y un algoritmo de pedido automático para prever demanda y generar pedidos diarios a la central de compras.
\item Reduje un 30 \% las caducidades y un 80 \% los productos sin venta durante más de un mes mediante reposición próxima al just-in-time, sin afectar a las ventas.
\item Programé un sistema de redistribución de stock entre tres tiendas según stock y rotación, con un segundo ciclo diario para productos centralizados y reposición el mismo día.
\item Reorganicé el trabajo en equipos polivalentes por tienda para que caja, reposición y secciones de frescos se apoyaran según la carga de trabajo.
\item Implanté Trello para asignar, registrar y seguir tareas operativas de responsables de tienda y de sección; revisaba las tareas que asignaba y ejecutaba las que tenía asignadas personalmente.
\item Realicé cuadres de caja durante mi primer año y diseñé posteriormente un sistema en Excel que mejoró el proceso final de cuadre de todas las cajas.
\end{itemize}

\section*{Trayectoria}
\textbf{Director Ejecutivo | Herfrailes S. L. | 01/10/2011–03/10/2024}

\section*{Formación}
Bachillerato y COU | IES Saulo Torón, Gáldar\\

\end{document}
